\scanpage{7}

Considérons le cas d'un \uline{automorphisme} $f$, et voyons quel est le \unsure{polynôme} \unsure{du} contragrédient
\[
\begin{aligned}
P_{\check f}(t)&=\deleted{\det}(1-\check f t)=\prod(1-\lambda_i^{-1}t)\\
&\deleted{=\unsure{\prod}\lambda_i^{-1}(\lambda_i-t)}\\
&=\prod_i(-\lambda_i^{-1}t)(1-\lambda_i t^{-1})\\
&=\frac{(-1)^n t^n}{\prod\lambda_i}\deleted{\unsure{\det}}\prod(1-\lambda_i t^{-1})\\
&=\frac{(-1)^n t^n}{\det f}P_f(\unsure{-t^{-1}|t^{-1}})
\end{aligned}
\]
donc
\deleted{
\[
P_{\check f}(t)=\deleted{\frac{\unsure{(-1)^n t^n}}{\det f}}\unsure{\varepsilon(f)}t^nP_f(t^{-1})\qquad\unsure{\varepsilon(f)}=
\]
}
\[
\boxed{P_{\check f}(t)=(-t)^{\varepsilon(f)}\det f^{-1}\,P_f(t^{-1})}
\]
i.e
\[
\boxed{L_{\check f}(t)=(-t)^{-\varepsilon(f)}\det f\,L_f(t^{-1})}
\]
Bien entendu, la dernière n'a \unsure{pas} de sens quand on \unsure{substitue|remplace} $\unsure{P_f^{\pm}(t)|P_f(t)}$ par une série formelle quelconque : il est essentiel de se borner aux séries formelles qui sont produits \unsure{alternés} finis de \uline{polynômes} [à terme constant $1$, à terme dominant inversible -- cette dernière condition n'étant utile que si $k$ est un anneau \unsure{général|quelconque} $\ldots$].

\mspage{6}
